% !TeX root = ../second_main.tex
\chapter*{Introduzione}
\addcontentsline{toc}{chapter}{Introduzione}

% --- §0 EPIGRAFE -------------------------------------------------------
\begin{flushright}
\itshape
«Un sistema pensionistico non è una macchina\\
per pagare assegni: è un meccanismo\\
per redistribuire il rischio tra le generazioni.»\\[4pt]
\upshape
Elsa Fornero, \textit{Chi ha paura delle riforme} (2018)
\end{flushright}

\vspace{1em}

% --- §1 AGGANCIO NARRATIVO ---------------------------------------------
\noindent
Il dibattito previdenziale italiano ha storicamente privilegiato due domande:
\textit{quando} si può andare in pensione e \textit{quanto} costerà farlo.
Entrambe sono legittime. Entrambe, però, presuppongono come risolto il problema
più profondo: \textbf{se} il sistema è in grado di garantire a chi lavora oggi
un'esistenza dignitosa nella vecchiaia. È da questa terza domanda che muove la
presente tesi.

La previdenza pubblica obbligatoria risponde a fallimenti di mercato documentati
dalla teoria economica: i mercati privati non riescono a offrire copertura completa
contro il rischio demografico aggregato, le asimmetrie informative distorcono
l'offerta assicurativa, e la miopia intertemporale degli individui genera risparmi
sistematicamente insufficienti in assenza di un obbligo contributivo
\cite{Artoni2014, Bosi2018}. Nel sistema a ripartizione adottato dall'Italia,
i lavoratori attivi finanziano le pensioni degli anziani coevi, nella fiducia che
le generazioni successive faranno altrettanto: un patto intergenerazionale che
funge da meccanismo di redistribuzione del rischio demografico collettivo
\cite{Bosi2018}.

Un sistema alternativo, quello a capitalizzazione, lega il rendimento ai mercati
finanziari anziché alla crescita del PIL. Il passaggio da un sistema all'altro 
comporta però un costo di transizione di estrema complessità per l'Italia, 
data una spesa pensionistica attestata al 15,4\% del PIL nel 2024 \cite{MEF_RGS2025}: 
la generazione che affronta il passaggio si troverebbe infatti a dover finanziare 
simultaneamente le pensioni già in erogazione e il proprio accumulo futuro. 
Il dibattito tra i due modelli, pur fondamentale sul piano teorico, si scontra 
quindi con ostacoli attuativi di notevole portata nel breve e medio periodo \cite{Artoni2014}.

La domanda che rimane aperta, e che quasi nessuna riforma ha mai affrontato
con rigore, riguarda l'adeguatezza delle prestazioni future. Il dibattito
previdenziale ha attraversato decenni di interventi che hanno modificato
\textit{come} si calcola la pensione (riforma Dini, 1995), \textit{quando}
ci si può ritirare (riforma Fornero, 2011), e \textit{quanto tempo prima}
è possibile farlo cedendo a pressioni di breve periodo (Quota~100 nel 2019,
Quota~103 nel 2023). Ciascuno di questi interventi ha risposto a una domanda
contingente senza tornare al problema fondativo: garantire a tutti, e non
soltanto ai lavoratori già prossimi alla quiescenza, una vecchiaia
economicamente sostenibile.

% --- §2 IL PARADOSSO CENTRALE -----------------------------------------
\section*{Il paradosso della sostenibilità senza adeguatezza}

Il sistema previdenziale italiano si trova in una condizione di
\textbf{sostenibilità fragile}. Dal punto di vista macroeconomico, la traiettoria
appare rassicurante: la spesa per pensioni si attesta al 15,4\% del PIL nel 2024
\cite{MEF_RGS2025}, valore elevato nel confronto europeo ma destinato a ridursi
strutturalmente dopo il picco previsto intorno al 2042-2044, quando l'ultima
coorte del \emph{baby boom} raggiungerà l'età pensionabile portando il rapporto
spesa/PIL vicino al 17\%, per poi scendere al di sotto del 14\% nel 2070 grazie
all'applicazione generalizzata del metodo contributivo e alla stabilizzazione del
rapporto demografico tra pensionati e occupati \cite{MEF_RGS2025}.

Questa traiettoria discendente è il prodotto di una scelta tecnica con un costo
sociale profondo. Il sistema contributivo nozionale (NDC,
\emph{Notional Defined Contribution}), introdotto dalla riforma Dini e
generalizzato dalla riforma Fornero, calcola la pensione come funzione dell'intera
vita lavorativa e la comprime sistematicamente all'aumentare della speranza di
vita, attraverso la revisione periodica dei coefficienti di trasformazione.
In formula, la pensione annua lorda è:

\begin{equation}
P = \frac{MC_{NDC}}{e_L}
\qquad \text{con} \qquad
MC_{NDC} = \sum_{j=1}^{L} C_j \,(1+g)^{L-j}
\label{eq:ndc}
\end{equation}

\noindent dove $MC_{NDC}$ è il montante contributivo nozionale,
$C_j = c \cdot S_j$ i contributi versati nell'anno $j$
(con $c$ aliquota contributiva e $S_j$ retribuzione),
$g$ il tasso di capitalizzazione nozionale pari alla media mobile quinquennale
del PIL nominale, $e_L$ la speranza di vita all'età di pensionamento $L$
\cite{Bosi2018, MEF_RGS2025}.

\paragraph{Applicazione numerica.}
Si consideri un lavoratore dipendente con 35 anni di contributi, retribuzione
lorda costante di 30.000 euro e aliquota $c = 33\%$ (contributi annui: 9.900 euro),
con tasso nozionale $g = 1{,}5\%$, assunto a fini esemplificativi come ipotesi
prudenziale di crescita del PIL nominale \cite{ISTATCapitalizzazione2025, ISTATCapitalizzazione2024, ISTATCapitalizzazione2023}.
\[
\begin{aligned}
MC_{NDC}
  &= 9.900 \cdot \frac{(1{,}015)^{35} - 1}{0{,}015} \\
  &\approx 9.900 \times 45{,}6 \\
  &\approx 451.440 \text{ euro}
\end{aligned}
\]

Con il coefficiente di trasformazione a 67 anni fissato al 5,608\% per il
biennio 2025-2026 \cite{MEF_RGS2025}, la pensione lorda annua è:

\[
P_{\text{lordo}} = 0{,}05608 \times 451.440 = 25.320 \text{ euro annui}
\]

Il confronto al lordo è però parziale: il trattamento fiscale e contributivo
di lavoratori e pensionati diverge in modo significativo. Il lavoratore versa
contributi a proprio carico pari al 9,19\% della retribuzione lorda \cite{INPSAliquote2023}(circa
2.757 euro), che riducono l'imponibile IRPEF; la sua aliquota media effettiva
complessiva si colloca attorno al 27,7\% sul lordo. 

Per capire in che misura la pensione riesca a sostituire il precedente reddito lavorativo, dobbiamo calcolare la pensione netta e confrontarla con l'ultimo stipendio netto. Il pensionato non versa contributi e gode di un regime fiscale disciplinato dall'art. 13 del Testo Unico delle Imposte sui Redditi (TUIR) \cite{TUIR917_1986}. Per un reddito di 25.320 euro, che rientra nello scaglione tra 8.500 e 28.000 euro (comma 3, lett. b), a cui si aggiunge la maggiorazione per i redditi tra 25.000 e 29.000 euro (comma 3-bis), la detrazione spettante nel 2025 è:

\[
D_{\text{pens}} = 700 + \left[ 1.255 \times \frac{28.000 - 25.320}{19.500} \right] + 50 \approx 922 \text{ euro}
\]

L'IRPEF lorda (23\% $\times$ 25.320 = 5.824 euro) si riduce così a 4.902 euro;
con le addizionali regionali e comunali (media nazionale 1,5\%, circa 380 euro), l'imposta totale è di 5.282 euro. L'aliquota media effettiva si attesta quindi al 20,9\%, per una pensione netta di:

\[
P_{\text{netto}} = 25.320 \times (1 - 0{,}209) = 20.038 \text{ euro annui}
= 1.541 \text{ euro mensili}
\]

Lo stipendio netto vale circa 21.700 euro annui (1.668 euro mensili). Il
\textbf{tasso di sostituzione netto} è quindi:

\[
TS_{\text{netto}} = \frac{20.038}{21.700} = 92{,}3\%
\]

contro l'84\% lordo. Il differenziale si spiega con l'assenza, per il
pensionato, dei contributi IVS (acronimo di Invalidità, Vecchiaia e Superstiti) che sul lavoratore equivalgono a un prelievo
implicito di quasi 10 punti percentuali. Questa asimmetria è rilevante nella
valutazione dell'adeguatezza: i confronti condotti esclusivamente su basi lorde
sottostimano il reddito disponibile del pensionato rispetto all'ultimo anno
di lavoro.

La correzione non attenua il problema demografico; lo inquadra meglio. Se il
coefficiente di trasformazione scende al 5,2\% nel 2040 per effetto dell'aumento
della speranza di vita, la pensione lorda cala a:

\[
P_{\text{lordo},2040} = 0{,}052 \times 451.440 = 23.475 \text{ euro}
\]

Applicando nuovamente l'art. 13 del TUIR su questo nuovo imponibile (la detrazione sale a circa 991 euro, senza maggiorazione perché il reddito è sotto i 25.000 euro), l'aliquota media effettiva scende al 20,3\%. La pensione netta diventa circa 18.715 euro, pari a 1.440 euro mensili. Il tasso di sostituzione netto scende così all'86,2\%: una perdita di 6,1 punti percentuali rispetto allo scenario base, conseguenza di un solo parametro su cui il lavoratore non esercita alcuna leva: la longevità aggregata della propria coorte.

Su carriere con discontinuità contributiva, tipiche della Generazione Z, le
proiezioni della Ragioneria Generale dello Stato portano il tasso di sostituzione
lordo dei dipendenti privati al 58,4\% nel 2070, dal 73,6\% del 2010
\cite{MEF_RGS2025}. Convertendo al netto, il guadagno fiscale è simmetrico tra
generazioni e non compensa la perdita strutturale di rendimento del sistema a
contribuzione definita. Un lavoratore entrato nel 1982 e pensionatosi nel 2020 ha
percepito un tasso di sostituzione lordo dell'81,5\%, pari a circa l'87\% al netto.
Un giovane entrato nel 2022 che andrà in pensione nel 2060 otterrà il 64,8\% lordo:
circa 17 punti percentuali in meno \cite{Censis_Confcooperative2026}. La previdenza
complementare risale al 66\% lordo \cite{MEF_RGS2025}, ma sull'ipotesi di adesione
piena che per i lavoratori discontinui rimane un'assunzione difficilmente
verificabile nella pratica.

La contraddizione è strutturale. \textbf{La sostenibilità finanziaria del sistema
è stata acquistata al prezzo dell'adeguatezza futura delle prestazioni.} Le riforme
hanno risolto il problema del come finanziare le pensioni di domani senza rispondere
al problema del se quelle pensioni saranno sufficienti a garantire una vecchiaia
dignitosa. Questo trade-off è il nucleo analitico attorno al quale si sviluppa
la presente tesi.

% --- §3 LE DOMANDE DI RICERCA -----------------------------------------
\section*{Le domande di ricerca}

La tesi si articola attorno a tre domande di ricerca strettamente interconnesse,
che procedono dal livello macroeconomico-istituzionale a quello individuale.

\begin{description}
\item[\textbf{Domanda 1 (adeguatezza).}]
Il sistema NDC italiano, nella sua configurazione attuale e nelle sue proiezioni di
lungo periodo, è in grado di garantire alle generazioni future un tenore di vita
dignitoso nella vecchiaia? In che misura il tasso di sostituzione atteso è compatibile
con le funzioni previdenziale, assicurativa e assistenziale che la teoria economica
assegna alla previdenza pubblica \cite{Bosi2018}?

\item[\textbf{Domanda 2 (riforme strutturali).}]
Quali interventi di politica economica, dai correttivi tecnici al sistema NDC
(revisione delle aliquote di capitalizzazione, garanzie di minimo pensionistico,
flessibilità attuarialmente neutra dell'età di uscita) al più radicale
\emph{Contributo sull'Automazione} come fonte alternativa di finanziamento,
possono correggere il trade-off tra sostenibilità e adeguatezza?

\item[\textbf{Domanda 3 (agentività individuale).}]
In quale misura il singolo lavoratore, in un sistema che ha progressivamente
trasferito su di lui il rischio demografico e il rischio di longevità, mantiene
la capacità di agire sulla propria condizione futura? E questa agentività è
equamente distribuita tra chi ha carriere regolari e chi le ha discontinue, o
costituisce essa stessa un ulteriore fattore di disuguaglianza?
\end{description}

% --- §4 STRUTTURA DELLA TESI ------------------------------------------
\section*{Struttura della tesi}

La tesi è organizzata in tre parti logicamente consequenziali, che corrispondono a
tre livelli di analisi progressivamente più orientati all'azione.

La \textbf{Prima Parte} costruisce le fondamenta teoriche, istituzionali e storiche
del sistema previdenziale italiano. Il primo capitolo ricostruisce la logica economica
della previdenza pubblica (fallimenti del mercato, funzione del patto
intergenerazionale, modelli a ripartizione e a capitalizzazione) per poi
descrivere il funzionamento del sistema NDC e calcolarne il tasso di rendimento
implicito, con riferimento alla letteratura canonica \cite{Bosi2018, Artoni2014}.
Una sezione di sintesi storica chiude il capitolo ripercorrendo l'evoluzione del
sistema pensionistico italiano dalla riforma Amato del 1992 alla riforma Fornero
del 2011, analizzando le deroghe di Quota 100, Quota 102 e Quota 103 come
manifestazioni di un conflitto irrisolto tra equità attuariale e consenso elettorale.

La \textbf{Seconda Parte} affronta il nucleo del problema, analizzando il trade-off
tra sostenibilità e adeguatezza. Il secondo capitolo funge da \textbf{diagnosi
empirica}: attraverso le proiezioni di lungo periodo della Ragioneria Generale dello
Stato \cite{MEF_RGS2025}, identifica il picco di spesa atteso e quantifica il
deterioramento del tasso di sostituzione per le generazioni future, confrontando
scenari con e senza previdenza complementare. Il terzo capitolo passa dalla diagnosi
alle \textbf{soluzioni}, mappando le proposte di riforma strutturale in discussione
nel dibattito internazionale, dalla letteratura di \textsc{VoxEU} ai contributi di
\textsc{Lavoce.info} e del \textsc{CERP}, valutandone l'efficacia e la fattibilità
politica.

La \textbf{Terza Parte} sposta il fuoco dall'analisi macroeconomica di sistema
all'azione del singolo, ponendo una domanda precisa: in un contesto di riforme
rinviate e rischio demografico crescente, quanto può davvero fare il lavoratore
per tutelarsi? Il quarto capitolo, attraverso simulazioni su diversi profili di
carriera della Generazione Z, esamina la \textbf{regressività implicita} della
previdenza complementare e delinea un \textit{framework} di pianificazione
individuale basato su risparmio e riscatti contributivi, condizionato però da
soglie minime di reddito.

Le \textbf{Conclusioni} raccolgono le risposte alle tre domande di ricerca e
formulano raccomandazioni di policy, distinguendo tra misure adottabili nel breve
periodo, senza necessità di riforma organica, e interventi strutturali di lungo
respiro.